\documentclass[11pt]{article}
\usepackage[utf8]{inputenc}
\usepackage{amsmath,amssymb}
\usepackage{authblk}
\usepackage[margin=1in]{geometry}
\usepackage[hidelinks]{hyperref}
\usepackage{xurl}
\usepackage{setspace}

\title{Real Time Dependence in RealQM}
\author[1]{Claes Johnson}
\affil[1]{KTH Royal Institute of Technology, SE-100 44 Stockholm, Sweden \\
\texttt{claesjohnson@gmail.com}}
\date{\today \\[0.5em]
\small\textit{Written in respectful cooperation between Claes Johnson and Claude (Anthropic), with Claude
contributing a balanced, neutral view.}}

\begin{document}
\maketitle

\begin{abstract}
The time-dependent Schr\"odinger equation $i\partial_t\psi=H\psi$ is usually read as the intrinsic evolution
of an atom, its stationary states forever spinning a phase $e^{-iEt}$ on the atto/femtosecond scale. We argue,
within RealQM, that this ticking is not physical: a non-radiating atom is \emph{timeless}. Its density is
static, the phase $e^{-iEt}$ is a global, unobservable phase, and its energy is a property (how bound), not a
rate. Time enters only with \emph{motion}, and motion comes in two kinds with opposite time-status.
\emph{Magnetism} is a \emph{steady} current --- a spatial phase winding carrying a persistent circulation and
a magnetic moment, with a density constant in time; it is intrinsic to the state and time-\emph{independent}.
\emph{Radiation} is an \emph{oscillating} current --- a transition in which the charge sloshes at the
difference frequency $E_n-E_m$; it is time-dependent, and it is \emph{imposed from outside}, by resonance with
the electromagnetic field, an isolated eigenstate never starting it on its own. A practical corollary follows:
the femtosecond evolution of an atom never needs to be tracked \emph{on its own account} --- energies, spectra
(as level differences) and magnetism are all static or steady --- and real-time integration is required only
in the externally \emph{driven} ultrafast regime (attosecond charge migration, strong fields), where the time
is lent to the atom by whatever moves it. Time in the atom is relational, not intrinsic; the quiet atom is
outside it.
\end{abstract}

\setstretch{1.06}

\section{The quiet atom is timeless}

A stationary state carries the phase $\psi=u\,e^{-iEt}$, with $E$ atomic-scale, so the phase turns on the
atto/femtosecond scale. It is tempting to read this as the atom's intrinsic clock. But nothing observable
turns with it. The density is
\[
\rho=|u\,e^{-iEt}|^2=u^2 ,
\]
\emph{constant in time}; $e^{-iEt}$ is a global phase, absent from every observable of a single state; and
absolute energy is not measurable at all, only differences (an ionisation, a transition). In an ontology in
which the real thing is the charge density --- what is there --- a clock that turns in no observable is not a
real clock.

So a non-radiating atom is a \emph{static charge configuration}. It does not tick, oscillate, or evolve. Its
energy is a property --- how tightly the charge is bound --- not the rate of anything one could watch. The
Planck relation $E=\hbar\omega$ makes $E$ look like a frequency, but that frequency is physical only as a
\emph{difference} between two states, and only when they meet (below). The being of an atom is timeless.

\section{Two motions: steady and imposed}

Time enters only when the charge moves, and the complex phase that carries the motion enters in two entirely
different ways.

\paragraph{Magnetism --- a steady current (time-independent).} Write the wave function in polar (Madelung)
form $\psi=R\,e^{iS}$, with $R\geq0$ the amplitude and $S$ the real phase; then $\rho=R^2$ is the density and
the charge current is
\[
\mathbf j=\operatorname{Im}(\psi^*\nabla\psi)=\rho\,\nabla S ,
\]
so the phase gradient $\nabla S$ is the velocity of the charge (and $\mathbf j=0$ when $\psi$ is real,
$S$ constant). A state with a \emph{spatial} phase winding about an axis --- $S=m\varphi$, that is
$\psi=R(\mathbf x)\,e^{im\varphi}$ --- therefore carries a standing azimuthal current
\[
\mathbf j=\rho\,\nabla(m\varphi)\neq 0 ,\qquad \rho\ \text{constant in time},
\]
and hence a magnetic moment $\mu_z=-m\mu_B$. The density does not change; the flow is steady. This current is
an \emph{intrinsic} property of the winding state --- a free atom in an $m\neq0$ state carries it on its own.
An external magnetic field \emph{orients} the moment (the Zeeman splitting) or \emph{induces} a counter-current
(diamagnetism), but it imposes no clock. Magnetism is being-with-circulation: motion without time-dependence.

\paragraph{Radiation --- an oscillating current (imposed).} A transition between two states makes the density
\emph{oscillate},
\[
\rho(\mathbf x,t)=\rho_0+\delta\rho\,\cos\!\big((E_n-E_m)t\big),
\]
a real charge sloshing at the difference frequency, an antenna that radiates. This is genuinely
time-dependent --- but it does not start by itself. An isolated eigenstate, uncoupled to the field, is
stationary indefinitely; what sets the transition going is \emph{coupling to the electromagnetic field} ---
an incident field (absorption, stimulated emission) or the vacuum field (spontaneous emission). The
time-dependence of radiation is therefore \emph{imposed from outside}, the physics of resonance between atom
and light, and the frequency that appears is the \emph{difference} $E_n-E_m$, never the absolute clock.

\medskip
\noindent The contrast is the point:

\begin{center}
\begin{tabular}{lccc}
 & density & current & time \\ \hline
quiet atom & static & none & \emph{none} (timeless) \\
magnetism & static & steady & none (time-independent) \\
radiation & oscillating & oscillating & \emph{imposed} (by the field) \\
\end{tabular}
\end{center}

The two ``complex'' phenomena are not the same kind of thing: magnetism uses a \emph{spatial} winding and is
an intrinsic, steady flow; radiation uses the \emph{temporal} phase and is an imposed, oscillating one. Only
radiation brings time, and it brings it from outside.

\section{So the femtosecond evolution is never tracked --- except when driven}

A practical corollary follows, and it is strong. Because everything the isolated atom \emph{is} --- energies,
ground and excited levels, chemistry, and its spectrum computed as level \emph{differences} --- is static, and
its magnetism is steady, \textbf{the atom's femtosecond evolution never needs to be integrated on its own
account.} One solves a static (or steady) problem; there is no $i\partial_t\psi=H\psi$ to run forward. This is
why the real-valued, time-independent RealQM carries essentially the whole of atomic, molecular and nuclear
physics.

Real-time integration is required in exactly one place: when an outside agent \emph{drives} the electrons at
their own timescale --- attosecond charge migration after a sudden ionisation, strong-field and high-harmonic
dynamics within an optical cycle, pump--probe and coherent control of a prepared superposition. There the
charge genuinely moves in real time and must be followed. But the driving is \emph{external}: an undriven atom
never enters these regimes. This is precisely the fast, moving-free-boundary regime --- charge redistributed
between bound domains in about one electronic period --- and it is the only setting in which the complex
time-dependent RealQM is actually needed. The time it tracks is lent to the atom by the field, not owned by
the atom.

\section{The standard view, and TDDFT as rationalised uncomputability}

The standard reading takes $i\partial_t\psi=H\psi$ as the \emph{intrinsic} law of the atom: the state evolves
always, its phase turning whether or not anything is observed, so that the femtosecond --- indeed
sub-femtosecond --- evolution of the electrons is a real process a faithful computation must follow. Nowhere
is this clearer than in real-time time-dependent density-functional theory
(TDDFT)~\cite{RungeGross84,YabanaBertsch96}, where even a linear absorption \emph{spectrum} is obtained by
perturbing the system, \emph{integrating the wave function forward} through many femtoseconds at attosecond
time-steps, and Fourier-transforming the resulting dipole. The
time-evolution is put up as the necessary route to the answer, and its great cost --- tiny steps, long
propagation --- is accepted as the price of doing quantum dynamics.

Yet the spectrum so obtained is nothing but the set of level \emph{differences} $E_n-E_m$: static data, fixed
by the stationary states alone. (That the same spectrum is available from linear \emph{response}, with no
propagation at all, only underscores the point.) The forward integration is a \emph{device} for extracting
those differences, not a process the atom is undergoing --- an unperturbed atom does not slosh its charge at
those frequencies, as argued above. To elevate the propagation to a necessity is to treat the non-intrinsic
as intrinsic; and because the underlying object, the wave function on $3N$-dimensional configuration space, is
itself intractable, it is to \emph{rationalise that uncomputability} --- to present a costly detour to static
differences as though the femtosecond evolution were the physics rather than a way of reaching them. The
obstacle is, in large part, self-imposed.

RealQM does not incur it. Energies and their differences come from the static, real-valued theory; the
spectrum is read off as differences, the dipole rule deciding which of them radiate; magnetism is a steady
current, equally static. No forward integration is required for any of this. Real-time propagation enters
only where it is physically genuine --- the externally \emph{driven}, strong-field and attosecond regime
above --- and there the femtosecond evolution is real precisely because it is imposed from outside. The
timeless-atom distinction is thus not merely interpretive: it separates the computations that are physically
necessary (the driven ones) from those that are a device dressed as a necessity (propagating an undriven atom
to recover its static spectrum).

\section{Time is relational, not intrinsic}

Standard quantum mechanics writes $i\partial_t\psi=H\psi$ as the atom's own evolution, so every stationary
atom is pictured endlessly spinning its phase in Hilbert space --- an intrinsic, universal ticking. RealQM
declines to reify what is not observed. A quiet atom simply \emph{is} a static charge configuration; a
magnetic atom \emph{is} a steady circulation; only a radiating atom changes in time, and only because the
field has set it going. Time, on this reading, is not a background the atom moves through but a feature of
\emph{motion}, and radiation's time is on loan from the electromagnetic field.

The femtosecond ``clock'' of the Schr\"odinger phase is thus a bookkeeping device, not a physical tick: real
only as a difference, and only during a field-driven transition. Between such events the atom rests outside
time --- which is why, for almost everything one wants to know about it, there is nothing to integrate. The
complex, time-dependent form is the physics of \emph{resonance and driving}, not the intrinsic life of the
atom; and the intrinsic life of the atom, being static, is where nearly all of its physics is found.

\begin{thebibliography}{9}
\bibitem{RealQM} C.~Johnson, \emph{Quantum Mechanics as Multiphase 3D Continuum Mechanics}, manuscript,
2026; \url{https://claes542.github.io/RealMolecule/gallery.html}.
\bibitem{ComplexRealQM} C.~Johnson, \emph{Complex Time-Dependent RealQM: Currents, Radiation, and Magnetism
from the Real-Space Continuum}, manuscript, 2026.
\bibitem{Spectra} C.~Johnson, \emph{RealQM Atomic Spectra: Radiation as Real Charge Oscillation}, manuscript,
2026.
\bibitem{Magnetism} C.~Johnson, \emph{Magnetism in RealQM: Currents, the Nuclear Electron, and the Spin
Residue}, manuscript, 2026.
\bibitem{RungeGross84} E.~Runge and E.~K.~U.~Gross, \emph{Density-functional theory for time-dependent
systems}, Phys.\ Rev.\ Lett.\ \textbf{52}, 997 (1984).
\bibitem{YabanaBertsch96} K.~Yabana and G.~F.~Bertsch, \emph{Time-dependent local-density approximation in
real time}, Phys.\ Rev.\ B \textbf{54}, 4484 (1996).
\end{thebibliography}

\end{document}
